% Part: first-order-logic
% Chapter: sequent-calculus
% Section: proof-theoretic-notions

\documentclass[../../../include/open-logic-section]{subfiles}

\begin{document}

\begin{editorial}
  આ વિભાગ સિક્વન્ટ કલન માટે સાબિત કરી શકાય તેવા હોવાના સંબંધ અને
  સુસંગતતાની વ્યાખ્યાઓ એકત્ર કરે છે.
\end{editorial}
\sourcecorrection{OLSEQ-003}{મૂળની
\texttt{proof-theoretic-notions.tex}ની સંપાદકીય નોંધ આ વ્યાખ્યાઓ
``સ્વાભાવિક નિગમન માટે'' હોવાનું કહે છે, પરંતુ ફાઇલ સિક્વન્ટ-કલન
પ્રકરણમાં છે અને તેની બધી વ્યાખ્યાઓ $\Log{LK}$-સિક્વન્ટો દ્વારા આપવામાં
આવી છે. તેથી અહીં ``સિક્વન્ટ કલન માટે'' લખ્યું છે.}

\iftag{FOL}
      {\olfileid{fol}{seq}{ptn}}
      {\olfileid{pl}{seq}{ptn}}

\olsection{સાબિતી-સૈદ્ધાંતિક ખ્યાલો}

\begin{explain}
આપણે અનેક મહત્વના અર્થાનુસારી ખ્યાલો (પ્રામાણ્ય, ફલિતતા અને સંતોષ્યતા)
વ્યાખ્યાયિત કર્યા છે; એ જ રીતે હવે તેમને અનુરૂપ \emph{સાબિતી-સૈદ્ધાંતિક
ખ્યાલો} વ્યાખ્યાયિત કરીએ છીએ. આ ખ્યાલો !!{structure}sમાં !!{sentence}sના
સંતોષનો આશરો લઈને નહિ, પરંતુ અમુક સિક્વન્ટોની !!{derivability} અથવા
!!{nonderivability}નો આશરો લઈને વ્યાખ્યાયિત થાય છે. આ બંને પ્રકારના ખ્યાલો
મળી જાય છે એ એક મહત્વની શોધ હતી. તેઓ મળે છે તે જ \emph{યથાર્થતા} અને
\emph{પૂર્ણતાના પ્રમેય}નો વિષય છે.
\end{explain}

\begin{defn}[પ્રમેયો]
જો સિક્વન્ટ $\quad \Sequent !A$ની $\Log{LK}$માં કોઈ !!a{derivation}
હોય, તો !!^a{sentence}~$!A$ \emph{પ્રમેય} છે. $!A$ પ્રમેય હોય તો
$\Proves !A$ અને ન હોય તો $\Proves/ !A$ લખીએ છીએ.
\end{defn}

\begin{defn}[!!^{derivability}]
!!^a{sentence}~$!A$ એ !!{sentence}sના ગણ~$\Gamma$માંથી
\emph{!!{derivable}} છે, $\Gamma \Proves !A$, જો અને માત્ર જો એવો સાન્ત
ઉપગણ~$\Gamma_0 \subseteq \Gamma$ અને $\Gamma_0$માંનાં !!{sentence}sની
એવી શ્રેણી $\Gamma_0'$ હોય કે $\Log{LK}$ એ
$\Gamma_0' \Sequent !A$ને !!{derive}s કરે. $!A$ એ $\Gamma$માંથી
!!{derivable} ન હોય તો $\Gamma \Proves/ !A$ લખીએ છીએ.
\end{defn}

સંકોચન, શિથિલીકરણ અને અદલાબદલીના નિયમોને કારણે $\Gamma_0'$માંનાં
!!{sentence}sનો ક્રમ અને તેમની સંખ્યા મહત્વ ધરાવતાં નથી: જો
$\Gamma_0' \Sequent !A$ !!{derivable} હોય, તો $\Gamma_0'$માં હોય એ જ
!!{sentence}s ધરાવતી કોઈપણ $\Gamma_0''$ માટે
$\Gamma_0'' \Sequent !A$ પણ નિષ્પન્ન થઈ શકે છે. દાખલા તરીકે, જો
$\Gamma_0 = \{!B, !C\}$ હોય, તો $\Gamma_0' = \tuple{!B, !B, !C}$ અને
$\Gamma_0'' = \tuple{!C, !C, !B}$ બંને એવી શ્રેણીઓ છે જેમાં માત્ર
$\Gamma_0$નાં !!{sentence}s જ છે. તેમાંથી એક ધરાવતું સિક્વન્ટ
!!{derivable} હોય, તો બીજું ધરાવતું સિક્વન્ટ પણ નિષ્પન્ન થઈ શકે છે; જેમ કે:
\begin{prooftree}
  \AxiomC{}
  \Deduce$!B, !B, !C \fCenter !A$
  \RightLabel{\LeftR{\Contraction}}
  \UnaryInf$!B, !C \fCenter !A$
  \RightLabel{\LeftR{\Exchange}}
  \UnaryInf$!C, !B \fCenter !A$
  \RightLabel{\LeftR{\Weakening}}
  \UnaryInf$!C, !C, !B \fCenter !A$
\end{prooftree}
હવેથી, $\Gamma_0$ એ !!{sentence}sનો સાન્ત ગણ હોય ત્યારે
$\Gamma_0 \Sequent !A$ એ એવું કોઈપણ સિક્વન્ટ છે જેનું પૂર્વાંગ
$\Gamma_0$માંનાં !!{sentence}sની શ્રેણી છે, એમ કહીશું; અને જરૂર પડે ત્યાં
સંકોચન, અદલાબદલી અને શિથિલીકરણોને ગર્ભિત રીતે સમાવિષ્ટ માનીશું.

\begin{defn}[સુસંગતતા]
વાક્યોનો ગણ~$\Gamma$ \emph{અસુસંગત} છે જો અને માત્ર જો એવો સાન્ત
ઉપગણ~$\Gamma_0 \subseteq \Gamma$ હોય કે $\Log{LK}$ એ
$\Gamma_0 \Sequent \quad$ને !!{derive}s કરે. $\Gamma$ અસુસંગત ન હોય,
એટલે કે દરેક સાન્ત $\Gamma_0 \subseteq \Gamma$ માટે $\Log{LK}$ એ
$\Gamma_0 \Sequent \quad$ને !!{derive} ન કરે, તો તેને \emph{સુસંગત}
કહીએ છીએ.
\end{defn}

\begin{prop}[સ્વવાચકતા]
\ollabel{prop:reflexivity}
જો $!A \in \Gamma$ હોય, તો $\Gamma \Proves !A$.
\end{prop}

\begin{proof}
આરંભિક સિક્વન્ટ $!A \Sequent !A$ !!{derivable} છે અને
$\{!A\} \subseteq \Gamma$ છે.
\end{proof}

\begin{prop}[એકદિશવર્ધિતા]
\ollabel{prop:monotonicity}
જો $\Gamma \subseteq \Delta$ અને $\Gamma \Proves !A$ હોય, તો
$\Delta \Proves !A$.
\end{prop}

\begin{proof}
માનો કે $\Gamma \Proves !A$, એટલે કે એવો સાન્ત
$\Gamma_0 \subseteq \Gamma$ છે કે $\Gamma_0 \Sequent !A$
!!{derivable} છે. $\Gamma \subseteq \Delta$ હોવાથી $\Gamma_0$ એ
$\Delta$નો પણ સાન્ત ઉપગણ છે. તેથી $\Gamma_0 \Sequent !A$ની
!!{derivation} એ $\Delta \Proves !A$ પણ બતાવે છે.
\end{proof}

\begin{prop}[પરંપરિતતા]
\ollabel{prop:transitivity}
જો $\Gamma \Proves !A$ અને $\{!A\} \cup \Delta \Proves !B$ હોય, તો
$\Gamma \cup \Delta \Proves !B$.
\end{prop}

\begin{proof}
જો $\Gamma \Proves !A$ હોય, તો કોઈ સાન્ત $\Gamma_0 \subseteq \Gamma$
અને $\Gamma_0 \Sequent !A$ની !!a{derivation}~$\pi_0$ હોય છે. જો
$\{!A\} \cup \Delta \Proves !B$ હોય, તો કોઈ સાન્ત
$\Delta_0 \subseteq \Delta$ માટે $!A, \Delta_0 \Sequent !B$ની
!!a{derivation}~$\pi_1$ હોય છે. નીચેની !!{derivation} જુઓ:
\begin{prooftree}
  \AxiomC{}
  \RightLabel{$\pi_0$}
  \Deduce$\Gamma_0 \fCenter !A$
  \AxiomC{}
  \RightLabel{$\pi_1$}
  \Deduce$!A, \Delta_0 \fCenter !B$
  \RightLabel{\Cut}
  \BinaryInf$\Gamma_0, \Delta_0 \fCenter !B$
\end{prooftree}
$\Gamma_0 \cup \Delta_0 \subseteq \Gamma \cup \Delta$ હોવાથી આ
$\Gamma \cup \Delta \Proves !B$ બતાવે છે.
\end{proof}

નોંધો કે ખાસ કરીને તેનો અર્થ એવો થાય છે કે જો $\Gamma \Proves !A$ અને
$!A \Proves !B$ હોય, તો $\Gamma \Proves !B$. વળી, જો
$!A_1, \dots, !A_n \Proves !B$ અને દરેક~$i$ માટે
$\Gamma \Proves !A_i$ હોય, તો $\Gamma \Proves !B$ એવું પણ ફલિત થાય છે.

\begin{prop}
\ollabel{prop:incons}
$\Gamma$ અસુસંગત હોય જો અને માત્ર જો દરેક વાક્ય~$!A$ માટે
$\Gamma \Proves {!A}$.
\end{prop}

\begin{proof}
અભ્યાસપ્રશ્ન.
\end{proof}

\begin{prob}
\olref[fol][seq][ptn]{prop:incons} સાબિત કરો.
\end{prob}

\begin{prop}[સઘનતા]
  \ollabel{prop:proves-compact}
  \begin{enumerate}
  \item જો $\Gamma \Proves !A$ હોય, તો એવો સાન્ત ઉપગણ
    $\Gamma_0 \subseteq \Gamma$ છે કે $\Gamma_0 \Proves !A$.
  \item જો $\Gamma$નો દરેક સાન્ત ઉપગણ સુસંગત હોય, તો $\Gamma$ સુસંગત છે.
  \end{enumerate}
\end{prop}

\begin{proof}
  \begin{enumerate}
    \item જો $\Gamma \Proves !A$ હોય, તો એવો સાન્ત ઉપગણ
      $\Gamma_0 \subseteq \Gamma$ છે કે સિક્વન્ટ
      $\Gamma_0 \Sequent !A$ને !!a{derivation} હોય. તેથી
      $\Gamma_0 \Proves !A$.
    \item જો $\Gamma$ અસુસંગત હોય, તો એવો સાન્ત ઉપગણ
      $\Gamma_0 \subseteq \Gamma$ છે કે $\Log{LK}$ એ
      $\Gamma_0 \Sequent \quad$ને !!{derive}s કરે. પરંતુ ત્યારે
      $\Gamma_0$ એ $\Gamma$નો અસુસંગત સાન્ત ઉપગણ છે.
  \end{enumerate}
\end{proof}

\end{document}
